\section{Introduction}
\label{sec:introduction}

Autonomous navigation of unmanned aerial vehicles (UAVs) in environments with limited GNSS signal availability is one of the central challenges of modern mobile robotics. Satellite signal degradation can occur in dense urban environments, under heavy forest canopy, in tunnels and canyons, and under deliberate jamming or spoofing conditions~\cite{yue2025airground, radwan2024uav, delmerico2018benchmark, qin2018vinsmono}. In such scenarios, a UAV must rely on alternative information sources: RGB and thermal cameras, inertial measurement units (IMU), altimeters, LiDARs, and reference maps~\cite{nieuwenhuisen2015autonomous, nieuwenhuisen2016autonomous, yue2024semantic, hu2025review}.

Contemporary autonomous navigation systems perform pose estimation, map construction, semantic segmentation, object recognition, occupancy forecasting, and trajectory planning as separate, loosely coupled modules~\cite{feng2025survey, placed2022survey}. This approach leads to computational redundancy -- the same images are repeatedly analyzed by different modules, and semantic and geometric information is not efficiently shared between perception and planning stages~\cite{rosen2024scene, ding2025survey}.

In this work, we propose S4D-TAM (Semantic 4D Token Attention Map), a novel navigation system architecture in which the fundamental environment representation is a persistent, semantic, continuously updated map of 4D tokens derived from the SLAM point cloud. RGB and thermal camera data are used primarily to update the geometry and semantic meaning of points and objects stored in the map. Planning, risk assessment, and decision-making modules do not repeatedly analyze the same images -- they operate on a compressed spatial representation containing geometry, semantics, dynamics, observation history, and uncertainty~\cite{pomianek2026s4dtam}.

The key scientific novelty of the proposed system is the introduction of a persistent memory map composed of semantic 4D tokens, which serves as a unified representation for localization, environment understanding, future scene state prediction, risk estimation, and UAV trajectory planning. The system remembers not only where an object was located, but also when and how many times it was observed, from which sensors the information originated, how reliable its position and recognition are, whether the object is static, dynamic, or periodically displaced, what spatial relations hold between it and other terrain elements, what invisible areas may be hidden behind it, and what significance the object has for flight safety and planning~\cite{pomianek2026s4dtam}.

The article combines a conceptual system architecture with an executable research reference and a reproducible benchmark framework. The current repository implements token state and lifecycle management, proposal and association, multimodal encoder interfaces, attention-related processing, uncertainty calibration utilities, reference-map and topological support, causal occupancy and motion forecasting, deterministic risk-aware planning, telemetry, and the integrated \texttt{S4DTAMReference} execution path. The final learned hierarchical model, full public-dataset reproduction, and confirmatory results remain under development. Section~\ref{sec:experiments} therefore specifies a prospective two-level validation protocol rather than reporting final performance results.
